\documentclass[12pt,a4paper]{article}

% --- Encoding & Language ---
\usepackage[utf8]{inputenc}
\usepackage[T1]{fontenc}
\usepackage[english]{babel}
\usepackage{microtype}

% --- Layout ---
\usepackage[top=3cm, bottom=3cm, left=3cm, right=3cm]{geometry}
\usepackage{parskip}
\setlength{\parskip}{8pt}

% --- Fonts ---
\usepackage{lmodern}

% --- Tables ---
\usepackage{booktabs}
\usepackage{longtable}
\usepackage{array}
\usepackage{tabularx}
\renewcommand{\arraystretch}{1.5}

% --- Math ---
\usepackage{amsmath}

% --- Graphics ---
\usepackage{graphicx}
\usepackage{float}
\usepackage{tikz}
\usetikzlibrary{arrows.meta, positioning, shapes.geometric}

% --- Code Listings ---
\usepackage{listings}
\lstset{
    basicstyle=\ttfamily\small,
    breaklines=true,
    frame=single,
    backgroundcolor=\color{white},
    numbers=left,
    numberstyle=\tiny,
    numbersep=8pt,
    xleftmargin=14pt,
    framexleftmargin=14pt,
    showstringspaces=false,
    captionpos=b,
    aboveskip=10pt,
    belowskip=6pt,
}
\usepackage{xcolor}

% --- Headers & Footers ---
\usepackage{fancyhdr}
\pagestyle{fancy}
\fancyhf{}
\fancyhead[L]{\small Elevator Control System -- Group E (C2)}
\fancyhead[R]{\small System Architecture -- HTWG Konstanz}
\fancyfoot[C]{\small\thepage}
\renewcommand{\headrulewidth}{0.4pt}
\renewcommand{\footrulewidth}{0pt}

% --- Section Styling ---
\usepackage{titlesec}
\titleformat{\section}
  {\large\bfseries}{\thesection}{1em}{}
  [\vspace{-4pt}\rule{\linewidth}{0.5pt}\vspace{4pt}]
\titleformat{\subsection}
  {\normalsize\bfseries}{\thesubsection}{1em}{}
\titleformat{\subsubsection}
  {\normalsize\bfseries}{\thesubsubsection}{1em}{}

\titlespacing{\section}{0pt}{24pt}{8pt}
\titlespacing{\subsection}{0pt}{16pt}{4pt}
\titlespacing{\subsubsection}{0pt}{12pt}{4pt}

% --- Hyperlinks (black) ---
\usepackage{hyperref}
\hypersetup{
    colorlinks=false,
    hidelinks,
    pdftitle={Elevator Control System -- Group E},
    pdfauthor={Ruben Kaiser, Lukas Rank, Yasin Maier, Mevlut Atilla},
}

% --- Misc ---
\usepackage{enumitem}
\setlist[itemize]{itemsep=4pt, topsep=6pt, leftmargin=1.5em}
\setlist[enumerate]{itemsep=4pt, topsep=6pt, leftmargin=1.5em}

% New section always on a new page
\newcommand{\sectionbreak}{\clearpage}

% ======================================================================
\begin{document}
% ======================================================================

% --- Title Page ---
\begin{titlepage}
\thispagestyle{empty}
\begin{center}

\vspace*{2cm}

\rule{\linewidth}{1.5pt}
\vspace{0.5cm}

{\Huge\bfseries Elevator Control System}\\[0.5cm]
{\Large\bfseries and Human Machine Interface}

\vspace{0.4cm}
\rule{\linewidth}{1.5pt}

\vspace{1.5cm}

{\large System Architecture -- Project Documentation}\\[0.3cm]
{\normalsize Faculty of Electrical Engineering and Information Technology}\\[0.2cm]
{\normalsize HTWG Konstanz -- University of Applied Sciences}

\vspace{2.5cm}

{\large\bfseries Group E (C2)}

\vspace{1cm}

\begin{tabular}{ll}
\toprule
\textbf{Name} & \textbf{Role} \\
\midrule
Ruben Kaiser      & Control System \\
Lukas Rank        & Control System \\
Yasin Maier       & Human Machine Interface \\
Mevl\"{u}t Atilla & Human Machine Interface \\
\bottomrule
\end{tabular}

\vfill

{\normalsize Summer Semester 2026}

\end{center}
\end{titlepage}

% --- TOC ---
\newpage
\thispagestyle{empty}
\tableofcontents
\newpage

% ======================================================================
\section{Project Overview}
% ======================================================================

\subsection{Group Members and Responsibilities}

Group E consists of four members organised into two sub-teams.
Each member is responsible for at least one project domain.

\begin{table}[H]
\centering
\begin{tabularx}{\textwidth}{llX}
\toprule
\textbf{Name} & \textbf{Team} & \textbf{Primary Domain} \\
\midrule
Ruben Kaiser      & Control System & Control logic, FSM, Modbus I/O mapping, overall architecture \\
Lukas Rank        & Control System & MQTT transport (control side), in-memory simulation, unit tests \\
Yasin Maier       & HMI            & Swing HMI, shaft animation, MQTT integration (HMI side) \\
Mevl\"{u}t Atilla & HMI           & Console HMI, message contract design, integration testing \\
\bottomrule
\end{tabularx}
\caption{Group members, teams and primary domains}
\end{table}

\subsection{Project Plan and Time Estimation}

The project ran over approximately three active weeks in June 2026.
The table below shows the estimated hours per person and domain.
The total of roughly 160 hours corresponds to the expected workload
for a 6-ECTS project (2 persons $\times$ 6 weeks $\times$ 12\,h/week).

\begin{table}[H]
\centering
\begin{tabular}{lrrrrr}
\toprule
\textbf{Domain} & \textbf{Ruben K.} & \textbf{Lukas R.} & \textbf{Yasin M.} & \textbf{Mevl. A.} & \textbf{Total} \\
\midrule
Architecture / Design   & 10\,h &  6\,h &  4\,h &  4\,h & 24\,h \\
Domain \& Control Logic & 18\,h &  8\,h &  2\,h &  2\,h & 30\,h \\
Modbus / Simulation     &  6\,h & 10\,h &  0\,h &  2\,h & 18\,h \\
MQTT Transport          &  4\,h & 10\,h &  8\,h &  6\,h & 28\,h \\
HMI (Swing / Console)   &  0\,h &  2\,h & 18\,h & 10\,h & 30\,h \\
Testing                 &  4\,h &  8\,h &  2\,h &  4\,h & 18\,h \\
Documentation           &  4\,h &  2\,h &  4\,h &  2\,h & 12\,h \\
\midrule
\textbf{Total} & \textbf{46\,h} & \textbf{46\,h} & \textbf{38\,h} & \textbf{30\,h} & \textbf{160\,h} \\
\bottomrule
\end{tabular}
\caption{Estimated development effort in hours per person and domain}
\end{table}

\subsection{Weekly Work Report}

\begin{longtable}{@{}lp{3.6cm}p{3.6cm}p{3.6cm}@{}}
\toprule
\textbf{Week} & \textbf{Ruben K. / Lukas R.} & \textbf{Yasin M.} & \textbf{Mevl\"{u}t A.} \\
\midrule
\endfirsthead
\toprule
\textbf{Week} & \textbf{Ruben K. / Lukas R.} & \textbf{Yasin M.} & \textbf{Mevl\"{u}t A.} \\
\midrule
\endhead
\midrule
\multicolumn{4}{r}{\small\itshape continued on next page} \\
\endfoot
\bottomrule
\endlastfoot

Week 1 \newline (16.06)
  & Architecture design, domain model, ports \& adapters scaffold, Level / Sensor enums
  & Initial HMI concept, MQTT contract discussion
  & MQTT contract discussion, console HMI skeleton \\[6pt]

Week 2 \newline (23.06)
  & Control FSM (SCAN logic), deceleration chain, RequestStore, SimulatedPlcGateway, 21 unit tests
  & Swing HMI: shaft animation, cabin rendering, call buttons, event log
  & Console HMI completion, MQTT command routing, integration testing \\[6pt]

Week 3 \newline (30.06)
  & MQTT wiring on control side, lab broker auth, deploy scripts, presentation
  & HMI broker auth, live status display, emergency / reset buttons, presentation
  & End-to-end testing with lab broker, documentation \\

\caption{Weekly work report by person}
\end{longtable}

% ======================================================================
\section{Design Considerations}
% ======================================================================

\subsection{Architectural Style: Hexagonal Architecture (Ports \& Adapters)}

The system follows a \textbf{hexagonal (ports \& adapters)} architecture.
The core control logic is completely isolated from I/O concerns -- it has no
knowledge of Modbus, MQTT, or any external system.

\begin{figure}[H]
\centering
\ttfamily\small
\begin{tabular}{ccccc}
  & \textit{(driving port)} & & \textit{(driven ports)} & \\[4pt]
  HMI / Tests
  & $\xrightarrow{\quad\textup{commands}\quad}$
  & \fbox{\begin{tabular}{c}Application\\[-2pt]\scriptsize ElevatorControlService\end{tabular}}
  & $\xrightarrow{\quad\textup{status/events}\quad}$
  & HMI sink \\[16pt]
  & & $\downarrow$ & & \\[4pt]
  & & \fbox{Domain (pure logic)} & & \\[16pt]
  & & $\downarrow$ & & \\[4pt]
  & & \fbox{PlcGateway} & $\longrightarrow$ & PLC / Sim \\
\end{tabular}
\rmfamily
\caption{Hexagonal architecture: the domain is isolated from all I/O by ports}
\end{figure}

This design was chosen because the assignment (§1.3) explicitly requires components
with well-defined, unidirectional interfaces. Each interface is a port; each concrete
I/O implementation is a replaceable adapter.

\subsection{Module Structure}

The project is split into two independent Maven modules under \texttt{Control-HMI/}:

\begin{itemize}
    \item \texttt{SysArch\_Control} -- the elevator control system
    \item \texttt{SysArch\_HMI} -- the reference HMI (graphical Swing + headless console)
\end{itemize}

The two modules share \textbf{only} the JSON message contract in package
\texttt{de.htwg.sysarch.mqtt} (a copy lives in each). There is no
cross-module Java import between control and HMI.

\subsubsection{Control System Package Map}

\begin{table}[H]
\centering
\begin{tabularx}{\textwidth}{lX}
\toprule
\textbf{Package / Class} & \textbf{Responsibility} \\
\midrule
\texttt{domain/model/}                      & Value objects: Level, Direction, Speed, Drive, LevelSensors, \ldots \\
\texttt{domain/control/ElevatorController}  & FSM -- one \texttt{step()} call per 100\,ms cycle \\
\texttt{domain/control/RequestStore}        & Pending cabin and hall calls; SCAN predicates \\
\texttt{application/ElevatorControlService} & Orchestrates one cycle; implements OperatorPanel \\
\texttt{application/ControlLoop}            & Fixed-rate 100\,ms scheduler \\
\texttt{application/port/in/OperatorPanel}  & Driving port: HMI sends calls / commands here \\
\texttt{application/port/out/PlcGateway}    & Driven port: read sensors / write actuators \\
\texttt{application/port/out/HmiGateway}    & Driven port: publish status \& events to HMI \\
\texttt{infrastructure/modbus/}             & Modbus/TCP adapter (ModbusIoMap, ModbusPlcGateway) \\
\texttt{infrastructure/simulation/}         & In-memory PLC simulation for offline development \\
\texttt{infrastructure/mqtt/}               & MQTT adapter (MqttHmiGateway, MqttCommandRouter) \\
\texttt{infrastructure/console/}            & Interactive stdin driver (\texttt{--interactive} mode) \\
\bottomrule
\end{tabularx}
\caption{Control system package map}
\end{table}

\clearpage
\subsection{Interface Convention}

Following §1.3 of the assignment, interface variables are \textbf{writable in one
direction only}:

\begin{itemize}
    \item \textbf{OperatorPanel} (driving) -- the HMI writes commands into the
          control system: cabin calls, hall calls, emergency stop, and reset.
    \item \textbf{HmiGateway} (driven) -- the control system writes state out to
          the HMI: immutable \texttt{ElevatorStatus} snapshots and an
          \texttt{ElevatorEvent} stream.
\end{itemize}

The HMI has no direct write access to any internal control state.

\subsection{Control Logic: SCAN Algorithm}

The elevator implements the \textbf{collective-control (SCAN / look) algorithm}:

\begin{enumerate}
    \item When idle, move toward the oldest pending request.
    \item While moving, stop at every request in the current direction
          (matching hall call or cabin call at that level).
    \item When no more requests exist ahead, reverse if requests exist behind.
    \item Otherwise become idle.
\end{enumerate}

\subsubsection{Finite State Machine}

The \texttt{ElevatorController} cycles through the following phases on each
100\,ms tick:

\begin{table}[H]
\centering
\begin{tabularx}{\textwidth}{lX}
\toprule
\textbf{Phase} & \textbf{Description} \\
\midrule
\texttt{IDLE}         & No pending requests; motor off. \\
\texttt{MOVING}       & Travelling at full speed (V2) toward the next target level. \\
\texttt{DECELERATING} & Approach sensor triggered; speed reduced to V1. \\
\texttt{CRAWLING}     & Safety sensor triggered; motor set to crawl for fine positioning. \\
\texttt{DWELL}        & Cabin at level: door opens, minimum 6\,s wait, door closes. \\
\texttt{EMERGENCY}    & Emergency stop engaged; motor off at any position. \\
\bottomrule
\end{tabularx}
\caption{Elevator FSM phases}
\end{table}

\clearpage
\subsubsection{Deceleration Chain}

When approaching a target level $T$, the motor steps down through three speed
stages triggered by the level sensors:

\begin{table}[H]
\centering
\begin{tabular}{llll}
\toprule
\textbf{Direction} & \textbf{Sensor} & \textbf{Offset} & \textbf{Action} \\
\midrule
UP   & approach lower (\texttt{al}) & $-1.5$\,m  & V2 $\rightarrow$ V1 \\
UP   & safety lower   (\texttt{sl}) & $-50$\,mm  & V1 $\rightarrow$ CRAWL \\
UP   & reached        (\texttt{r})  & $0$\,mm    & CRAWL $\rightarrow$ STOP \\
\midrule
DOWN & approach upper (\texttt{au}) & $+1.5$\,m  & V2 $\rightarrow$ V1 \\
DOWN & safety upper   (\texttt{su}) & $+50$\,mm  & V1 $\rightarrow$ CRAWL \\
DOWN & reached        (\texttt{r})  & $0$\,mm    & CRAWL $\rightarrow$ STOP \\
\bottomrule
\end{tabular}
\caption{Deceleration chain triggered by level sensors}
\end{table}

\subsection{HMI Interface Boundary}

The control system and the HMI run as \textbf{separate processes} and communicate
exclusively over \textbf{MQTT} (Eclipse Mosquitto on \texttt{ea-pc165:1883}).
They share only the JSON message contract -- neither process imports the other's
classes.

\begin{table}[H]
\centering
\begin{tabularx}{\textwidth}{llXl}
\toprule
\textbf{Topic} & \textbf{Direction} & \textbf{Payload} & \textbf{Retained} \\
\midrule
\texttt{/SysArch/E/status} & Control $\to$ HMI & \texttt{StatusMessage} JSON (full state snapshot) & yes \\
\texttt{/SysArch/E/event}  & Control $\to$ HMI & \texttt{EventMessage} JSON (arrival, door, alarm)  & no  \\
\texttt{/SysArch/E/cmd}    & HMI $\to$ Control & \texttt{CommandMessage} JSON (call, emergency, reset) & no \\
\bottomrule
\end{tabularx}
\caption{MQTT topic contract between control system and HMI}
\end{table}

\clearpage
Example payloads:

\begin{lstlisting}[language={}, caption={Status message (control to HMI)},
  basicstyle=\ttfamily\footnotesize, xleftmargin=0pt, framexleftmargin=0pt,
  numbers=none]
{
  "level": 2,  "direction": "UP",  "phase": "MOVING",
  "door": "CLOSED",  "velocity": 100,
  "emergencyActive": false,  "motorError": false,
  "cabinCalls": [4],  "hallUpCalls": [2],  "hallDownCalls": []
}
\end{lstlisting}

\begin{lstlisting}[language={}, caption={Command messages (HMI to control)},
  basicstyle=\ttfamily\footnotesize, xleftmargin=0pt, framexleftmargin=0pt,
  numbers=none]
{"type": "cabin",     "level": 3}
{"type": "hall",      "level": 2, "direction": "UP"}
{"type": "emergency", "action": "engage"}
{"type": "reset"}
\end{lstlisting}

% ======================================================================
\section{Test Plan and Test Execution}
% ======================================================================

\subsection{Test Strategy}

All tests are written with \textbf{JUnit 5}. Because the domain is pure Java
with no I/O dependencies, every FSM transition and SCAN predicate can be tested
without a PLC or MQTT broker. The MQTT adapters are verified against an
in-memory \texttt{MqttConnection} stub -- no running broker is required.

\subsection{Test Cases}

\begin{longtable}{@{}p{5cm}p{7cm}p{1.8cm}@{}}
\toprule
\textbf{Test method} & \textbf{What is verified} & \textbf{Result} \\
\midrule
\endfirsthead
\toprule
\textbf{Test method} & \textbf{What is verified} & \textbf{Result} \\
\midrule
\endhead
\endfoot
\bottomrule
\endlastfoot

\texttt{fullStopSequence\-ForCabinCallAbove}
  & FSM transitions IDLE $\to$ MOVING $\to$ DECELERATING $\to$ CRAWLING $\to$ DWELL
    on a cabin call above the current level.
  & PASS \\[4pt]

\texttt{hallCallPickupOn\-PassingFloor}
  & Elevator already moving up picks up a hall-up call at an intermediate floor.
  & PASS \\[4pt]

\texttt{directionReversal\-WhenNoRequestsAhead}
  & After serving all upward requests, controller reverses to DOWN when
    downward requests exist (req.~10).
  & PASS \\[4pt]

\texttt{emergencyStop\-DuringTravel}
  & Emergency stop halts the motor immediately; after reset, the elevator
    resumes toward the next scheduled stop (req.~12).
  & PASS \\[4pt]

\texttt{requestStoreScan\-Predicates}
  & \texttt{shouldStopAt()} and \texttt{hasRequestBeyond()} return correct
    results for mixed cabin and hall calls in both directions.
  & PASS \\[4pt]

\texttt{mqttHmiGateway\-PublishesStatus}
  & \texttt{MqttHmiGateway} serialises \texttt{ElevatorStatus} to the correct
    topic with well-formed JSON.
  & PASS \\[4pt]

\texttt{mqttCommandRouter\-DispatchesCabinCall}
  & JSON \texttt{\{"type":"cabin","level":3\}} on the cmd topic is routed
    to \texttt{OperatorPanel.cabinCall(LEVEL\_3)}.
  & PASS \\[4pt]

\texttt{mqttCommandRouter\-DispatchesEmergency}
  & \texttt{\{"type":"emergency","action":"engage"\}} triggers
    \texttt{OperatorPanel.emergencyStop()}.
  & PASS \\

\caption{Selected unit test results -- all 21 tests pass}
\end{longtable}

\subsection{Integration Test (Lab)}

The full system was verified end-to-end against the HTWG lab infrastructure:

\begin{itemize}
    \item Control system connected and authenticated to the MQTT broker at
          \texttt{ea-pc165:1883} using the group credentials for Group E.
    \item Swing HMI connected from a laptop via an SSH tunnel
          (\texttt{localhost:1883 $\to$ ea-pc165:1883}).
    \item Cabin and hall calls issued from the HMI were received and executed
          by the control system; the live shaft view updated accordingly.
    \item Emergency stop and reset round-trip verified.
\end{itemize}

\textit{Note:} The Modbus connection to the real PLC (\texttt{ea-pc111:506})
was not fully tested because the \texttt{easymodbus-maven} library integration
is still pending. All functional verification used the in-memory simulation mode.

% ======================================================================
\section{Development Effort}
% ======================================================================

The table below summarises the actual development effort based on the weekly
work reports. The project totals approximately 160 hours across all four members.

\begin{table}[H]
\centering
\begin{tabular}{lrrrrr}
\toprule
\textbf{Domain} & \textbf{Ruben K.} & \textbf{Lukas R.} & \textbf{Yasin M.} & \textbf{Mevl. A.} & \textbf{Total} \\
\midrule
Architecture / Design   & 10\,h &  6\,h &  4\,h &  4\,h & 24\,h \\
Domain \& Control Logic & 18\,h &  8\,h &  2\,h &  2\,h & 30\,h \\
Modbus / Simulation     &  6\,h & 10\,h &  0\,h &  2\,h & 18\,h \\
MQTT Transport          &  4\,h & 10\,h &  8\,h &  6\,h & 28\,h \\
HMI (Swing / Console)   &  0\,h &  2\,h & 18\,h & 10\,h & 30\,h \\
Testing                 &  4\,h &  8\,h &  2\,h &  4\,h & 18\,h \\
Documentation           &  4\,h &  2\,h &  4\,h &  2\,h & 12\,h \\
\midrule
\textbf{Total} & \textbf{46\,h} & \textbf{46\,h} & \textbf{38\,h} & \textbf{30\,h} & \textbf{160\,h} \\
\bottomrule
\end{tabular}
\caption{Development effort per person and domain}
\end{table}

% ======================================================================
\section{Project Evaluation and Potential Improvements}
% ======================================================================

\subsection{What Went Well}

\begin{itemize}
    \item \textbf{Architecture.}
          The hexagonal approach paid off immediately: the full SCAN control logic
          was implemented and tested against the in-memory simulation before any
          MQTT or Modbus code existed. Replacing the simulation with a real adapter
          required zero changes to the domain.

    \item \textbf{MQTT contract.}
          Agreeing on the shared JSON contract in week 1 allowed both sub-teams to
          develop independently. Integration was essentially friction-free.

    \item \textbf{Resilient transport.}
          The control system never depends on broker availability. Without
          \texttt{--mqtt}, or when the broker is unreachable, it falls back to
          console logging and keeps running. This made offline development
          straightforward.

    \item \textbf{Test coverage.}
          21 unit tests cover all FSM transitions, SCAN predicates, MQTT
          serialisation, and command routing without any external dependencies.
          Tests can be run with a single \texttt{mvn test}, with no lab access
          or running broker.
\end{itemize}

\subsection{Challenges}

\begin{itemize}
    \item \textbf{Modbus library integration.}
          The \texttt{easymodbus-maven} library coordinates were not confirmed
          in time. The address map is fully specified in \texttt{ModbusIoMap}
          and the adapter skeleton exists in \texttt{ModbusPlcGateway}, but no
          end-to-end PLC test was carried out.

    \item \textbf{Git Bash path mangling.}
          MQTT topics beginning with \texttt{/} are silently rewritten by MSYS
          path conversion in Git Bash. PowerShell helper scripts were added to
          work around this, but identifying the root cause cost significant
          debugging time.

    \item \textbf{Lab access constraints.}
          HTWG VPN or a physical lab seat is required for any live testing,
          which limited the time available for integration work.
\end{itemize}

\subsection{Potential Improvements}

\begin{enumerate}
    \item \textbf{Real Modbus connection.}
          The highest-priority open item is completing \texttt{ModbusConnection}
          with \texttt{easymodbus-maven} and verifying the I/O address map
          against the live PLC at \texttt{ea-pc111:506}.

    \item \textbf{User access levels in the HMI.}
          The assignment (§1.7 req.~5) requests two roles: Supervisor and
          Operator. A simple login dialog restricting the emergency and reset
          buttons to supervisors would fulfil this requirement.

    \item \textbf{Fairness statistics.}
          The optional person-simulation feature (§1.7 req.~7) was not
          implemented. It would require a per-person state machine
          (waiting $\to$ in cabin $\to$ arrived) and a statistics view showing
          minimum, maximum, and average travel time.

    \item \textbf{Shared contract as a Maven artifact.}
          The MQTT contract classes are currently duplicated in both modules.
          Publishing them as a versioned Maven artifact would eliminate the
          duplication and make breaking changes visible immediately.

    \item \textbf{One-step deployment.}
          The deploy script automates build and SCP to \texttt{ea-pc165}, but
          the SSH tunnel on the laptop must still be opened manually. A small
          companion script would make the full lab setup a single command.
\end{enumerate}

% ======================================================================
\appendix
\section{Requirements Traceability Matrix}
% ======================================================================

\begin{table}[H]
\centering
\begin{tabularx}{\textwidth}{clX}
\toprule
\textbf{\#} & \textbf{Requirement (§1.6)} & \textbf{Implementation} \\
\midrule
1  & 4 levels
   & \texttt{Level} enum: \texttt{LEVEL\_1} -- \texttt{LEVEL\_4} \\
2  & 2 speeds up/down + crawl + off
   & \texttt{Speed}, \texttt{Drive}, \texttt{MotorCommand}; mapped to Modbus coils and the crawl register in \texttt{ModbusIoMap} \\
3  & 2 approach sensors per level / direction
   & \texttt{LevelSensors}: approachLower, safetyLower, reached, safetyUpper, approachUpper \\
4  & Sensor for safe door position
   & \texttt{reached} sensor (0\,mm offset) \\
5  & Drive at highest speed, slow before stop
   & Deceleration: V2 $\to$ V1 $\to$ CRAWL $\to$ STOP in \texttt{ElevatorController} \\
6  & Hall call with direction; stop if same direction
   & \texttt{HallCall}, \texttt{RequestStore.shouldStopAt()} \\
7  & Cabin call selects target level
   & \texttt{RequestStore} cabin-call set \\
8  & Idle $\to$ oldest request
   & Timestamps in \texttt{RequestStore}; idle handling in \texttt{ElevatorController} \\
9  & Continue while requests ahead
   & \texttt{RequestStore.hasRequestBeyond()}, SCAN main loop \\
10 & Reverse when no more requests ahead
   & \texttt{ElevatorController.decideDirection()} \\
11 & Minimum dwell 6\,s + door cycle
   & Dwell timer in \texttt{ElevatorController} (DWELL phase) \\
12 & Emergency stop; resume after reset
   & EMERGENCY phase + \texttt{OperatorPanel.emergencyStop()} / \texttt{reset()} \\
\bottomrule
\end{tabularx}
\caption{Requirements traceability matrix (§1.6)}
\end{table}

% ======================================================================
\section{Configuration Reference}
% ======================================================================

\begin{table}[H]
\centering
\begin{tabularx}{\textwidth}{lX}
\toprule
\textbf{Parameter} & \textbf{Value (Group E / C2)} \\
\midrule
Modbus slave      & \texttt{ea-pc111.ei.htwg-konstanz.de:506} \\
MQTT broker       & \texttt{ea-pc165.ei.htwg-konstanz.de:1883} \\
MQTT base topic   & \texttt{/SysArch/E} \\
MQTT username     & \texttt{E} \quad (password not committed to version control) \\
Control cycle     & 100\,ms \\
Number of levels  & 4 \\
Level height      & 3.5\,m \\
Speed V1          & $\approx 0.1$\,m/s \\
Speed V2          & $\approx 1.0$\,m/s \\
Crawl speed range & $-5 \ldots +5$\,cm/s \\
\bottomrule
\end{tabularx}
\caption{System configuration for Group E}
\end{table}

\end{document}
